\clearpage\begingroup\SourceLeaf{60}\NoSourceHead
\begin{center}
{\Large [MAGNETISMUS UND GALVANISMUS.]}\\[1.5em]
{\large [AMTLICHER BERICHT.]}\par
\end{center}
\noindent\hspace*{3em}An
\begin{center}Königliches Universitäts-Curatorium.\end{center}
\begin{center}\textls[80]{Vortrag des Hofraths \textsc{Gauss} in Göttingen, das Bedürfniss eines
besonderen Locals für magnetische Beobachtungen betreffend.}\end{center}
\noindent\hspace*{\parindent}Unter den Gegenständen, welchen die Naturforscher des gegenwärtigen
Jahrhunderts das lebhafteste Interesse widmen, nimmt die Lehre vom Magnetismus,
wenn nicht die erste, doch eine der ersten Stellen ein, und seit den
letzten zwölf Jahren ist durch die grossen Entdeckungen von \textsc{Oersted}, \textsc{Ampère},
\textsc{Arago} und \textsc{Faraday} in Beziehung auf das wunderbare Band, welches
jene Lehre mit der von der Electricität und dem Galvanismus verknüpft,
jenes Interesse noch viel mehr gesteigert worden. Fast noch wichtiger aber,
als der glänzende Zuwachs unerwarteter Thatsachen, die in diesen Gebieten
entdeckt sind, ist der Umstand, dass auch hier die Versuche einer alles frühere
weit überflügelnden Schärfe, und ihre einfachen Grundgesetze einer wahrhaft
mathematischen Präcision fähig werden, so dass die Scheidewand zwischen
eigentlich sogenannter Physik und angewandter Mathematik auch hier (wie
längst in der Bewegungslehre und Optik) zu sinken, und die tiefer eingreifende
Bearbeitung dem Mathematiker anheim zu fallen anfängt.

Am meisten ist das letztere in Beziehung auf die den Erdmagnetismus
angehenden Beobachtungen der Fall, worüber die feinsten Beobachtungen
\clearpage\SourceLeaf{61}\HistoricalHead{56}{L}{PHYSIK, AMTLICHER BERICHT.}
\noindent eigentlich nur von solchen Mathematikern, die mit den feinsten Beobachtungsmitteln
vertraut sind, also von den praktischen Astronomen, erwartet werden
können. Immer kann jedoch der Einzelne nur durch seinen Standpunkt beschränktes
liefern, und eine kräftige Förderung unserer Einsichten in diesem
höchst interessanten Felde kann nur durch die vereinten Bemühungen vieler,
an \textls[120]{vielen Punkten} unseres Erdkörpers, erreicht werden. Hr. \textsc{von Humboldt}
hat hiezu bereits einen grossartigen Impuls gegeben, und eine Vereinbarung
zu begründen angefangen, welcher zum Besten der Wissenschaft die
ausgedehnteste Verbreitung zu wünschen ist. Er hat auf eigne Kosten in
Berlin ein besonderes Gebäude für magnetische Beobachtungen errichtet, in
welchem beim Bau durchaus gar kein Eisen gebraucht ist, und dasselbe mit
kostbaren von \textsc{Gambey} in Paris verfertigten Instrumenten ausgerüstet: hier
werden fortwährende Beobachtungen, besonders über die magnetische Declination
und ihre höchst merkwürdige, man kann sagen räthselhafte, tägliche
und stündliche Aenderung angestellt. Allein diese Beobachtungen gewinnen
erst dadurch ihre grösste Wichtigkeit, dass Hr. \textsc{von Humboldt} seine ausgebreiteten
Verbindungen benutzt hat und noch benutzt, um ähnliche correspondirend
gleichzeitige Beobachtungen mit eben solchen Hilfsmitteln an einer
Menge anderer Oerter an den entlegensten Punkten der Erde zu veranlassen.
Gegenwärtig finden solche Statt in Freiberg in Sachsen, in Paris, an vielen
Punkten des russischen Reichs, z. B. Petersburg, Moskwa, Nicolajew; selbst
in Peking und Südamerika, und schon jetzt sind mehrere interessante Resultate
aus dieser Vereinbarung hervorgegangen.

Vor ungefähr einem Jahre machte Hr. \textsc{von Humboldt} in einem Schreiben
an unsern Prof. \textsc{Weber} hierüber einige Mittheilungen, mit der Aufforderung,
auch Göttingen an diesem Plan Theil nehmen zu lassen, dessen Curatorium
stets gern, was den Wissenschaften erspriesslich ist, befördert. Es waren jenem
Briefe zugleich die Preise der Instrumente, wie \textsc{Gambey} in Paris sie liefert
und wie sie an den andern genannten Oertern angewandt werden, beigefügt:
diese Preise beliefen sich auf etwa 4000 Francs. Als Prof. \textsc{Weber} mir damals
diesen Brief mittheilte, konnte ich zwar nicht verkennen, wie sehr wünschenswerth
es sei, dass auch Göttingen an diesen rühmlichen Bestrebungen Theil
nehme, musste jedoch für rathsam halten, einen Antrag darauf einstweilen
noch zu verschieben, da ich selbst eben damals mit Versuchen beschäftigt
\clearpage\SourceLeaf{62}\HistoricalHead{57}{R}{ÜBER EIN LOKAL FÜR MAGNETISCHE BEOBACHTUNGEN.}\HistoricalFoot{XI1.}{8}
\noindent war, die magnetischen Apparate in ihrer ersten Grundlage wesentlich zu vervollkommnen.

Diese meine Bemühungen, bei welchen ich von dem Prof. \textsc{Weber} vielfältig
unterstützt bin, habe ich im Laufe des Jahres 1832 zur Reife zu bringen
gesucht: meine Erwartungen sind nicht bloss erfüllt, sondern weit übertroffen,
und ich habe die eingerichteten Apparate bereits zu einer Arbeit über die
Intensität des Erdmagnetismus angewandt, die den Gegenstand einer von mir
in der königl. Societät am 15. Dec. v. J. gehaltenen Vorlesung ausmacht[\textsuperscript{*})]\begingroup\renewcommand{\thefootnote}{}\footnotetext{[\textsuperscript{*}) \textit{Intensitas vis magneticae terrestris ad mensuram absolutam revocata}, vorgelegt am 15. December 1833, Werke V, S. 79; vergl. die Anzeige ebenda, S. 293.]}\endgroup.
Einen Bericht über diese Vorlesung in den hiesigen gelehrten Anzeigen\begingroup\renewcommand{\thefootnote}{**)}\footnote{Wie ich soeben erfahre, hat Hr. \textsc{von Humboldt} dieser Anzeige die Ehre erzeigt, sie zu schnellerer Verbreitung in Frankreich selbst ins Französische zu übersetzen.}\endgroup\ beehre
ich mich, beizulegen, da derselbe sowohl einen allgemeinen Begriff von
den neuen Apparaten und deren Leistungen zu geben dienen kann, als auch
die Umstände andeutet, die in dem gegenwärtigen --- wie in \textls[120]{jedem andern}
Local, welches nicht besonders für ähnliche Zwecke angelegt ist --- unmöglich
machen, zu ganz reinen und der Schärfe der Hülfsmittel würdigen Resultaten
zu gelangen.

Ich glaube nunmehro, meine Arbeiten zu derjenigen Reife gebracht zu
haben, um bei dem Königl. Universitäts-Curatorium auf Errichtung eines besonderen
Gebäudes für die verschiedenen auf den Erdmagnetismus sich beziehenden
Beobachtungen antragen zu können. Ich würde eine Pflicht gegen
die hiesige gelehrte Anstalt versäumen, wenn ich nicht jetzt auf die Gelegenheit
aufmerksam machte, wo Göttingen nicht allein auf eine rühmliche Art
in die Reihe anderer schon bestehender und sich immer mehr vervielfältigender
ähnlicher Anlagen eintreten, sondern durch ganz neue Methoden und Hülfsmittel
den andern vorleuchten kann.

In einem solchen Local würden zunächst und hauptsächlich diejenigen
Beobachtungen zu machen sein, für welche die von Hrn. \textsc{von Humboldt} errichteten
und veranlassten Etablissements dienen, also

1. die currenten jeden Tag einigemahl zu bestimmten Stunden anzustellenden
Beobachtungen der Variation der magnetischen Abweichung. Diese
würde in der Regel Prof. \textsc{Harding} gern auf sich nehmen.
\clearpage\SourceLeaf{63}\HistoricalHead{58}{L}{PHYSIK, AMTLICHER BERICHT.}\HistoricalFoot{}{}
\noindent\hspace*{\parindent}2. Ahnliche Beobachtungen an bestimmten, unter allen Beobachtern in
den verschiedenen Welttheilen verabredeten Tagen (acht im Jahre) von Viertelstunde
zu Viertelstunde, oder wenigstens von Stunde zu Stunde. Da diese
jedesmahl zwei Tage und eine Nacht en suite befassen, so sind sie nur durch
eine Vertheilung unter mehrere Personen ausführbar; ausser Prof. \textsc{Harding}
und mir wird Prof. \textsc{Weber} stets gern Theil daran nehmen, auch wird dabei
immer auf die Beihülfe dieses und jenes talentvollen Studirenden gerechnet
werden können.

Ausserdem würde ich selbst von Zeit zu Zeit die erst jetzt möglich
werdende viel delicatere Bestimmung der absoluten magnetischen Intensität
vornehmen.

Endlich würde diese Einrichtung zu fast unzähligen andern magnetischen,
galvano-magnetischen und elektromagnetischen Beobachtungen und Messungen
dienen, und gewiss manche bisher noch dunkle Partie in diesen Lehren dadurch
aufgehellt werden können.

Die von mir eingerichteten Apparate empfehlen sich übrigens auch dadurch,
dass sie vergleichungsweise nur geringe Kosten verursachen. Der oben
angegebene Preis der von Herrn \textsc{von Humboldt} und an den andern Orten
gebraucht werdenden Instrumente fällt mithin hier ganz weg, und es handelt
sich bloss um das Gebäude. Zwar würde eines besondern, ganz zu diesem
Zweck eingerichteten Locals nur ein Apparat von beträchtlich grössern Dimensionen,
als die bisher von mir gebrauchten, ganz würdig sein; allein auch
Apparate von solchen beträchtlich grössern Dimensionen würden nur geringe
Kosten machen, die hier gar nicht in Anschlag zu kommen brauchen, sondern
füglich unter den currenten kleinen Ausgaben für die Sternwarte verrechnet
werden könnten. Von den bisher schon vorhandenen beiden Apparaten
könnte etwa der eine demnächst an das physikalische Cabinet abgegeben
werden, während der andere für allerlei erforderliche comparative Messungen
vor wie nach seine Dienste leisten würde.\quad .\ .\ .\ .\hspace{1.1em}.\ .\ .\ .\ .\ .\ .\ .\ .\ .\ .\ .\ .\ .\ .

Göttingen, den 29. Januar 1833.
\begin{center}\rule{.18\textwidth}{.35pt}\end{center}
\clearpage\endgroup
